% Authoritative LaTeX manuscript; Markdown is a generated compatibility projection.

% <!-- part-title: Modeling -->

% <!-- chapter-title: Decision Models: What Should Happen Next? -->

\chapter{Decision Models: What Should Happen Next?}

A behavioral model can tell us that a failed job may be retried, escalated, or stopped. A decision model addresses the next question: given what we know about this failure, which alternative should we choose? A transient error may justify another attempt, while an exhausted retry budget may call for escalation. What makes this a decision model is that it represents the criteria for choosing among alternatives. The most useful form depends on the choice we face and the information available when we must make it.

\section{Conditions, Alternatives, and Decision Analysis}\label{heading:conditions-alternatives-and-decision-analysis}

We will again begin with a small example, examine the properties and analyses it supports, and then turn to DocAble. Here we will consider several kinds of representation, because choosing among alternatives can involve rules, permissions, constraints, or optimization. We begin with a decision table and introduce other forms as the questions require.

\subsection{Constructing a Decision Model}\label{heading:constructing-a-decision-model}

A service sells prepaid credits and also offers unlimited subscriptions. A customer requests work costing five credits. An incorrect admission decision could reject work the customer has paid to receive, or admit work without the required entitlement.

\textbf{Question.} Should the service admit the request, and if so, should it charge five credits? To answer, we need to know which entitlement applies. For a prepaid customer we also need the balance. For an unlimited subscriber we need to know whether the work is covered. \autoref{simple-admission-policy} assumes the request falls under exactly one entitlement.

{\footnotesize
\begin{longtable}[]{@{}
  >{\raggedright\arraybackslash}p{(\linewidth - 4\tabcolsep) * \real{0.3333}}
  >{\raggedright\arraybackslash}p{(\linewidth - 4\tabcolsep) * \real{0.3333}}
  >{\raggedright\arraybackslash}p{(\linewidth - 4\tabcolsep) * \real{0.3333}}@{}}
\caption[Decision modeling---admission by entitlement]{\emph{Decision modeling---admission by entitlement.} Each request falls under one entitlement. Unlimited entitlement admits covered work without drawing down finite credit; a finite balance must cover the requested work.}\label{simple-admission-policy}\tabularnewline
\toprule\noalign{}
\begin{minipage}[b]{\linewidth}\raggedright
Entitlement
\end{minipage} & \begin{minipage}[b]{\linewidth}\raggedright
Requested work
\end{minipage} & \begin{minipage}[b]{\linewidth}\raggedright
Decision
\end{minipage} \\
\midrule\noalign{}
\endhead
\bottomrule\noalign{}
\endlastfoot
Unlimited & Any amount covered by the subscription & Admit without drawing down finite credit \\
Finite balance & Enough credit for the request & Admit and charge the required credit \\
Finite balance & Insufficient credit for the request & Reject for insufficient credit \\
\end{longtable}
}

\textbf{Semantics.} Each row specifies conditions and the corresponding admission and charging decision. Read across the row whose conditions hold. A prepaid customer with five credits can proceed and spends the five; with four credits, the same request is rejected. For an unlimited subscription covering the work, the service admits the request without drawing on a balance.

To make the right decision, we must know the facts that could change it. Suppose a caller supplies only the number of prepaid credits. A customer who has spent the last credit and an unlimited subscriber who never bought credits can both have a balance of zero. The table gives them different answers. No revision to a rule that sees only the balance can recover the missing entitlement.

% <!-- point: permission-has-many-canonical-forms-unified-by-mood-not-shape | Decision tables, rule systems, and constraint systems make choices explicit; policy is one kind of decision model, expressed through access-control matrices, policy graphs, and feature models as well as those general forms. -->

A decision tree arranges the same choice as a sequence of tests: first determine the entitlement, then consult the balance only on the prepaid branch. Rules can express exceptions and their precedence; constraints can describe which combinations of features may be selected together. For permissions, an access-control matrix records which operations a subject may perform on an object~\citep{lampson1974}. When the alternatives differ in cost or performance, a quantitative model may help choose among them. Predicting a deployment's latency supplies one input; selecting the least expensive deployment that meets a latency target uses that input to make a decision.

% <!-- point: same-arrow-different-mood-is-the-sharpest-case-of-topology-not-fixing-meaning | An observed-call edge and a permitted-call edge can look identical while answering different questions; in the illustrated models, a missing edge means no recorded call or no permitted call respectively. -->

A permission graph is useful when the decision concerns a relationship. In \autoref{observed-vs-declared}, a web interface may call a report service, and the report service may call a renderer. To decide whether Web may call Renderer directly, look for that particular edge in the declared policy. The two-hop route through the report service grants no such permission. The observed graph records a direct Web-to-Renderer call. Comparing the two identifies a violation. If the observation had missed the call, its absence from the trace would tell us only what was recorded.

\begin{figure}[htbp]
\centering
\includegraphics[width=\linewidth]{book/latex/assets/observed-vs-declared-4186370076.pdf}
\caption{\emph{Observation compared with declared permission.} The observed web-to-render call violates the declared policy.}
\label{observed-vs-declared}
\end{figure}

% <!-- point: image-purpose-and-uncertainty-govern-the-decorative-decision | An image's informational role depends on context; DocAble uses a VLM classification with an explicit conservative policy that retains content when the classification is uncertain. -->

% Evidence: E03 | source: book/data/local-evidence/evidentiary-basis-11202026/integration/records/E03.md | kind: illustration
Some decisions cannot be made from the object alone. Accessibility standards distinguish informative images, which need a text alternative, from decorative images, which assistive technology can ignore~\citep{w3c2018wcag21}. Recall that a blue square might decorate a restaurant menu but be the subject of a lecture on modern art.\magefootnote{blue-monochrome-context}{For a real example, consider Yves Klein's \emph{Blue Monochrome} (1961), in the Museum of Modern Art's collection~\citep{moma-blue-monochrome}.} What matters is the image's contribution to the document, not merely what it depicts. A classifier can supply evidence about that role. The decision rule must say when the evidence warrants omitting a description and what to do when the role is uncertain.

\subsection{Properties and Analyses}\label{heading:decision-properties-and-analyses}

Decision models let us examine whether the available rules cover the cases they must handle, whether those rules agree, and whether the choices they permit satisfy the applicable constraints. When several acceptable choices remain, we can also ask which best serves a stated objective.

% <!-- point: a-declared-policy-turns-is-this-allowed-into-an-evaluation | Once permission is declared, the claims a policy designer wants become expressible — permitted/forbidden, mutual exclusion, prerequisites, consistency, completeness, least privilege — and "is this allowed?" becomes an evaluation rather than an argument. -->

A decision table is \textbf{complete} if it prescribes an outcome for every case in its declared scope. To check \autoref{simple-admission-policy}, examine the balance below, exactly at, and above the requested amount. If the finite rows used only strict inequalities, a customer with exactly five credits requesting five credits would match neither row. \textbf{Consistency} asks whether applicable rules can prescribe incompatible outcomes. If both rows included equality, that customer would match both the admitting and rejecting rows unless a precedence rule settled the conflict. A request covered by several entitlements presents a different issue: the table assumes one. Extending the policy to such requests requires deciding how the entitlements combine.

For a quantitative decision, first check \textbf{admissibility}: does an alternative satisfy every applicable constraint? Suppose a deployment must fit in 8~GB of memory and respond within 200~ms. A candidate requiring 9~GB is out, however fast or inexpensive it may be. Of the candidates that meet both limits, which costs least? That is an \textbf{optimality} question. A constraint solver can search for a feasible assignment; ranking feasible candidates by cost also requires an objective and comparable cost estimates. The least expensive candidate among those examined is an optimum over that set. A claim about all possible deployments needs a correspondingly wider analysis.

Checking \textbf{conformance} asks whether a particular decision or observed action agrees with the declared policy. For the service graph, compare the observed edges with the permitted set. For the entitlement table, use the inputs to select a row and compare its verdict with the service's actual response. Neither comparison settles whether the policy expresses the right judgment. A consistent policy could permit the very call the architect meant to exclude, just as an optimizer could minimize cost while omitting a requirement users care about. Engineers must examine the choice of inputs and criteria as well as the decisions made from them.

We now turn to DocAble to see how decision models address choices involving uncertain evidence, service permissions, and user entitlements.

\section{DocAble: Policies and Decisions}\label{heading:docable-policies-and-decisions}

DocAble uses decision models for several purposes: deciding whether an image needs a description, declaring which services may communicate, and determining whether a user may upload a document. The representations differ because the decisions depend on different information. Each must preserve the distinctions that can change its answer.

\subsection{When is an image decorative?}\label{heading:when-is-an-image-decorative}

% Evidence: E09 | source: book/data/local-evidence/evidentiary-basis-11202026/integration/records/E09.md | kind: dated-case-evidence
Accessibility remediation must distinguish images that convey information from those that are merely decorative. As \hyperref[2.5-decision:heading:conditions-alternatives-and-decision-analysis]{the worked example} illustrated, the same image can serve either purpose depending on its context. DocAble supplies a rendered slide to a vision-language model (VLM), which classifies each image or shape as clearly decorative, uncertain, or informational. The classification supplies evidence; the decision policy determines what to do with it. Only a clearly decorative verdict permits suppressing a description. An uncertain or missing verdict preserves the image's informational role. This asymmetry is deliberate: an unnecessary description is audible, while wrongly hiding information can leave a screen-reader user unaware that anything was omitted.

% Evidence: E09 | source: book/data/local-evidence/evidentiary-basis-11202026/integration/records/E09.md | kind: dated-case-evidence
The need for an explicit uncertainty policy emerged from evaluation. Earlier classifiers had marked content-bearing images as decorative, particularly in visualization-heavy documents. We tightened the policy by representing uncertainty explicitly and making its conservative interpretation common to the classification paths. The decision model now states how uncertain evidence affects remediation without prescribing which VLM must supply that evidence. A different VLM can make the semantic judgment; it does not redefine the policy for acting on it.

The remaining two DocAble examples concern permission. The system accepts documents for accessibility remediation, sends work among internal services, and returns corrected artifacts. Two decisions arise at different boundaries: which service may call another, and whether a user\textquotesingle s entitlement permits a requested upload. Both can be modeled as permission, but they need different representations.

\subsection{Which services may communicate?}\label{heading:which-services-may-communicate}

% Evidence: E10 | source: book/data/local-evidence/evidentiary-basis-11202026/integration/records/E10.md | kind: implementation
DocAble distributes document remediation across several services. The public web tier receives requests, while internal services perform document editing, generation, rendering, OCR, and other work. A typed service-flow model declares the permitted connections and is checked against the dependency map consumed by deployment. Each directed edge declares that one service may invoke another; an absent edge excludes that relationship from the declared policy.

\textbf{Service-flow model · Decision}

\begin{itemize}
\tightlist
\item
  \textbf{Engineering question} --- Which service may call which, and what may be reached from outside?
\end{itemize}

\begin{itemize}
\tightlist
\item
% Evidence: E10 | source: book/data/local-evidence/evidentiary-basis-11202026/integration/records/E10.md | kind: implementation
  \textbf{Model} --- a declared graph of services and the permitted edges between them, checked for correspondence with the dependency map consumed by deployment.
\end{itemize}

\begin{itemize}
\tightlist
\item
  \textbf{Property} --- internal communication is permitted only along a declared edge; one declared node is the public entry.
\end{itemize}

\begin{itemize}
\tightlist
\item
  \textbf{Quality attribute} --- least privilege; architectural integrity; security.
\end{itemize}

% Evidence: E10 | source: book/data/local-evidence/evidentiary-basis-11202026/integration/records/E10.md | kind: implementation
% <!-- point: the-service-flow-model-is-a-declared-source-of-truth-not-a-firewall | The typed graph declares permitted inter-service edges and is checked against the dependency map consumed by deployment; completeness checks concern the declaration, while runtime permission depends on installed controls. | terms: soft-vs-hard-governance -->

\autoref{service-flow-model} shows a reduced view of the service-invocation policy. The web tier is the only publicly reachable service. Its declared outgoing edges permit calls to selected internal services. Separately, the worker may invoke generation, rendering, and OCR services. The arrows express permissions, not observed traffic or a requirement that every permitted call occur. Asynchronous job dispatch is omitted from this selected view; absence here does not imply prohibition by the full policy.

% Evidence: E10 | source: book/data/local-evidence/evidentiary-basis-11202026/integration/records/E10.md | kind: implementation
The declared graph answers whether a proposed service invocation is permitted: look for the corresponding directed edge. It also supports reachability questions, such as which services the public web tier can reach through permitted calls. DocAble checks this typed graph against the dependency map consumed by deployment. A completeness check requires each gated edge to contain its declared connection and authentication information. These checks establish properties of the declarations and their correspondence, not whether runtime traffic follows the intended policy.

\begin{figure}[htbp]
\centering
\includegraphics[width=\linewidth]{book/latex/assets/service-flow-model-7634fa9071.pdf}
% Evidence: E10 | source: book/data/local-evidence/evidentiary-basis-11202026/integration/records/E10.md | kind: implementation
\caption{\emph{DocAble's service-flow model.} Arrows show selected permitted service invocations; asynchronous job dispatch is omitted. The web tier is the sole public entry point. The typed graph is checked against the deployment dependency map, while runtime conformance requires separate enforcement.}
\label{service-flow-model}
\end{figure}

% Evidence: E10 | source: book/data/local-evidence/evidentiary-basis-11202026/integration/records/E10.md | kind: implementation
% <!-- point: docable-does-not-reconcile-observed-traffic-so-correspondence-is-the-open-question | DocAble checks the typed graph against the deployment dependency map but does not reconcile observed traffic against it at runtime; whether the deployed services stay on the permitted edges, and what a departure should cost, is Alignment's question. -->

% Evidence: E10 | source: book/data/local-evidence/evidentiary-basis-11202026/integration/records/E10.md | kind: implementation
The declared graph can also serve as the policy against which observed service calls are checked. DocAble does not currently reconcile runtime traffic against that declaration. Whether to add such a check, and what consequences a violation should have, are Alignment decisions taken up in \magevalue{chapter:alignment}.

\subsection{Which entitlement authorizes an upload?}\label{heading:which-entitlement-authorizes-an-upload}

% <!-- point: the-entitlement-model-was-born-from-a-production-failure | The service-flow policy was declared in advance; DocAble's upload-admission decision acquired its model in the opposite order, after a production failure exposed a representation that could no longer express the decision. -->

% Evidence: E11 | source: book/data/local-evidence/evidentiary-basis-11202026/integration/records/E11.md | kind: dated-case-evidence
The service-flow model determines which calls are permitted. Upload admission presents a different decision: whether a user's entitlements authorize a particular document. This became a modeling problem when DocAble added unlimited-use groups to a system originally designed around finite credits.

% <!-- point: upload-entitlement-flows-through-groups | A user receives upload entitlement through groups — finite-credit or unlimited — so the admission decision depends on two facts: how many credits remain and whether any entitlement is unlimited. -->

\autoref{simple-admission-policy} assumed that exactly one entitlement applied to a request. DocAble's domain is more complex: a user may belong to several groups, each with its own entitlement. The upload decision must account for the available finite credits and whether any applicable group grants unlimited use (\autoref{entitlement-domain}). The production failure arose when these facts traveled separately through the system.

\begin{figure}[htbp]
\centering
\includegraphics[width=\linewidth]{book/latex/assets/entitlement-domain-14ff87d0d6.pdf}
% Evidence: E11 | source: book/data/local-evidence/evidentiary-basis-11202026/integration/records/E11.md | kind: implementation
\caption{\emph{The upload-entitlement domain.} A finite balance alone cannot determine admission: the model must also preserve whether any applicable group grants unlimited use.}
\label{entitlement-domain}
\end{figure}

% <!-- point: the-scalar-was-adequate-until-the-domain-changed | The original credit comparison was understandable directly in code; unlimited entitlement made the policy harder to follow across implementation sites and motivated a separate decision model. -->

% Evidence: E11 | source: book/data/local-evidence/evidentiary-basis-11202026/integration/records/E11.md | kind: dated-case-evidence
The original client upload rule was easy to understand from the code: if the user's credits covered the upload cost, allow the upload; otherwise, reject it. Box (1) in \autoref{entitlement-model-evolution} shows this balance-based check. A separate decision model would have added little to understanding that simple rule.

% Evidence: E11 | source: book/data/local-evidence/evidentiary-basis-11202026/integration/records/E11.md | kind: dated-case-evidence
Unlimited groups changed what the upload decision had to account for. As box (2) shows, a group could authorize an upload without contributing finite credits to the balance. The server's balance calculation, the client's check before sending an upload, and the server admission gate represented entitlement in different ways. Understanding and updating the intended policy now required checking how each part of the implementation handled both finite and unlimited entitlement.

% <!-- point: the-old-reduction-observed-the-distinction-and-discarded-it | The finite-credit reduction excluded unlimited groups, while the client upload check ignored a separate unlimited-status signal and rejected the upload before the correct server gate was reached. -->

% Evidence: E11 | source: book/data/local-evidence/evidentiary-basis-11202026/integration/records/E11.md | kind: dated-case-evidence
Box (3) traces the production failure. In August 2026, a user whose personal group provided unlimited use was blocked from uploading a thirty-eight-page presentation: \emph{the upload needs 38 credits, and you have 0}. The legacy server function read each group's unlimited flag twice: once to exclude unlimited groups from the credit sum, and once to compute a display class. For this user it returned zero credits and the class \emph{individual}. The login response omitted the detailed group records but carried a separate unlimited-status flag. The client upload check ignored that flag, compared the zero balance with the required 38 credits, and rejected the upload before sending the file. The server's admission logic, which handled this entitlement correctly, was never reached. The failure came from using the finite-credit total as an admission decision even though other system representations retained the unlimited entitlement.

% <!-- point: the-production-failure-produced-the-model | After fixing the client to read the existing unlimited signal, the team made the more complex entitlement policy explicit in a decision model and checked it with property tests, shared policy cases, and a shape check. -->

% Evidence: E11 | source: book/data/local-evidence/evidentiary-basis-11202026/integration/records/E11.md | kind: dated-case-evidence
The immediate fix made the client consult the unlimited-status flag it had been ignoring. That repaired the observed failure, but left the admission policy represented separately across the client and server. We therefore introduced the entitlement decision model shown in box (4): given the user's group entitlements and the requested work, should the upload be admitted, and under which entitlement? Property tests checked the model, a cross-runtime parity corpus checked shared policy cases in the server and client, and another check validated the model's shape. The display code derived its values from the model, while server admission delegated the policy decision to it.

\textbf{Entitlement decision model · Decision}

\begin{itemize}
\tightlist
\item
  \textbf{Engineering question} --- May this user upload this document, and under which entitlement?
\end{itemize}

\begin{itemize}
\tightlist
\item
  \textbf{Model} --- per-group entitlements that keep finite and unlimited as distinct cases; the output is a structured admission decision: admitted or not, the group that authorizes it, and the ordered plan by which the verdict was reached.
\end{itemize}

\begin{itemize}
\tightlist
\item
  \textbf{Property} --- an uncapped unlimited group can authorize an upload with zero finite credits; the display projection does not decide admission.
\end{itemize}

\begin{itemize}
\tightlist
\item
  \textbf{Quality attribute} --- decision correctness; fidelity between the domain and its representations.
\end{itemize}

% <!-- point: the-repair-was-model-to-model | The revised model yields a structured admission decision and a separate display projection; unlimited use produces an admitting decision and a display sentinel, even though the group has zero finite credits. -->

% Evidence: E11 | source: book/data/local-evidence/evidentiary-basis-11202026/integration/records/E11.md | kind: dated-case-evidence
The repaired model preserves the distinction that the earlier admission check lost. For the same uncapped unlimited group and 38-page request, it admits the upload and records the authorizing group and an \emph{admit\_unlimited} plan step, even though the finite credit balance remains zero. The display projection answers a different question: it returns a special numeric value for unlimited use, or a finite credit sum when all groups are finite. That display value cannot determine admission. The model can therefore be tested against cases involving finite credits, unlimited entitlement, and their combinations without treating a display value as the decision itself.

\begin{figure}[htbp]
\centering
\includegraphics[width=\linewidth]{book/latex/assets/entitlement-model-evolution-31ef2274cf.pdf}
\caption{\emph{The admission check must preserve unlimited entitlement.} Panels (1)–(4) explain the old representation, domain change, failure, and repair; they are not one runtime pipeline. In (3), the client uses the finite-credit total and ignores a separate unlimited signal. In (4), the same uncapped entitlement produces an admitting decision and a separate display value for unlimited use.}
\label{entitlement-model-evolution}
\end{figure}

% <!-- point: an-explicit-model-would-have-made-the-obligation-visible | The author changed the domain in one place and did not carry the change into every representation of the same decision; when the domain contains a consequential distinction, every representation used to reason about that decision must be capable of expressing it. -->

% Evidence: E11 | source: book/data/local-evidence/evidentiary-basis-11202026/integration/records/E11.md | kind: dated-case-evidence
A note in the first person: I changed the domain in one place --- groups could now be unlimited --- and did not carry the change into every representation of the same decision. An explicit model would not have made the change automatically correct; it would have made the inconsistency visible as an engineering obligation. When the domain contains a consequential distinction, every representation used to reason about that decision must be capable of expressing it. \magevalue{chapter:alignment} gives the recurring pattern here --- a failure converted into explicit structure that later work inherits --- a name.

The three decision models preserve different kinds of information: evidence about an image's role, permitted relationships between services, and the entitlements that authorize an upload. Each makes a decision inspectable by separating the information used to make it from the criteria governing the outcome. That separation also makes it possible to revise a classifier, service topology, or entitlement scheme without silently redefining the decision policy.

The next section examines how different models can be connected to answer questions that no individual representation can settle.
